\scanpage{30}
\mspage{29}
\[
\begin{aligned}
 \overset{\deleted{\unsure{\sim|\approx}}}{\unsure{M|\mathcal M}}^{*}
  &=\bigcup_{\deleted{\unsure{F}}\Phi\in\deleted{\unsure{\Phi}}\unsure{\mathcal F|\mathfrak F|\mathscr F}_M}
       \deleted{\unsure{F}}\Phi\subset\unsure{\mathcal P|P|\mathfrak P|\mathscr P}(\deleted{\unsure{F}}M)\\
 &\quad\unsure{\shortparallel|\text{''}}\\
 &\left\{\unsure{A|\dot A}\in\mathcal P(M)\mid\exists\,\deleted{\unsure{F}}\Phi\in\deleted{\unsure{\Phi}}\mathcal F_M
       \ \text{\unsure{avec}}\ A\in\deleted{\unsure{F}}\Phi\right\}
\end{aligned}
\]
i.e. l'ensemble des strates de toutes les figures $\in\deleted{\unsure{\Phi}}.\mathcal F_M$

\uline{At$^{*}$1}$_{\unsure{0|o}}$\quad $\overset{\deleted{\sim}}{M}^{*}$ est une figure ensembliste
\margin{\rule{0.45pt}{3.5\baselineskip}}
dans $M$, de support $M$. De plus, le préordre associé de $M$ \inserted{(\unsure{soit|sur} $\ll$)} est une relation d'\uline{ordre} i.e.
\[
\begin{aligned}
 \unsure{x|X},\unsure{y|Y}\in M,\ \unsure{x|X}\neq\unsure{y|Y}\quad\Longrightarrow\quad \exists\, A\in\overset{\deleted{\sim}}{M}^{*}
 &\quad\text{\unsure{avec}}\quad X\in A,\ Y\notin A\\
 &\quad\text{ou}\qquad Y\in A,\ x\notin A
\end{aligned}
\]

Ainsi, l'application
\[
 \unsure{x|X}\longmapsto \deleted{\unsure{M^{*}_{\ll x}|\widetilde M_{\ll x}}}M_{\ll\unsure{x|X}}
 \overset{\mathrm{def}}{=}\mathrm{Omb}(X)
\]
\[
 \underset{\unsure{\cdot\,\cdot}}{M}\longrightarrow\mathcal P(M)\tag*{(*)}
\]
induit une \uline{bijection} de $M$ \unsure{sur|avec} $\overset{\deleted{\sim}}{M}^{*}$, transformant $\ll$ en la relation d'inclusion.
Les $A\in\overset{\deleted{\sim}}{M}^{*}$ sont les parties fermées \uline{irréductibles} de $M$ admettant un pt générique.

\deleted{\uline{\unsure{At|Al}$^{*}$}2}\quad Soit $\overset{\deleted{\sim}}{\deleted{\unsure{F}}\Phi}\in\deleted{\unsure{\Phi}}\inserted{\mathcal F_M}$, [\unsure{donc|\deleted{donc}} $\overset{\deleted{\sim}}{F}\subset\mathcal P(M)$ et par définition même $\deleted{\unsure{F}}\Phi\subset\overset{\deleted{\sim}}{M}^{*}\subset\mathcal P(M)$), donc $\Phi$ est
